% !TEX root = main.tex
\section{Nghiên cứu liên quan}\label{sec:related}

\subsection{Các phương pháp truyền thống}\label{sec:related:classic}

\textbf{Đường đi ngắn nhất với trọng số tĩnh.} Thuật toán Dijkstra~\cite{dijkstra1959} giải bài toán đường đi ngắn nhất
một nguồn trên đồ thị có trọng số không âm bằng chiến lược tham lam: luôn ``chốt'' nút có nhãn khoảng cách nhỏ nhất
trong hàng đợi ưu tiên. Với hàng đợi nhị phân, độ phức tạp là $O((n+m)\log n)$. Thuật toán A*~\cite{hart1968} bổ sung một
hàm heuristic $h(v)$ ước lượng khoảng cách còn lại tới đích; nếu $h$ \emph{chấp nhận được} (không vượt quá khoảng cách
thật) thì A* vẫn tối ưu, và nếu $h$ \emph{nhất quán} thì mỗi nút được mở rộng tối đa một lần. Heuristic càng sát khoảng
cách thật, số nút phải mở rộng càng ít.

\textbf{Tăng tốc bằng tiền xử lý.} Trên mạng đường cỡ quốc gia, người ta tiền xử lý đồ thị để truy vấn nhanh hơn nhiều
bậc. ALT~\cite{goldberg2005alt} chọn một số \emph{điểm mốc} (landmark), tính sẵn khoảng cách tới/từ các điểm mốc và dùng
bất đẳng thức tam giác để xây heuristic chặt cho A*. Contraction Hierarchies~\cite{geisberger2008ch} ``co'' lần lượt các
nút ít quan trọng và thêm cạnh tắt, cho phép tìm kiếm hai chiều chỉ đi ``lên'' trong thứ bậc. Tổng quan đầy đủ về các
kỹ thuật này có trong~\cite{bast2016}. Thuật toán của Yen~\cite{yen1971} liệt kê $k$ đường đi không lặp ngắn nhất, thường
dùng để đưa ra tuyến thay thế.

\textbf{Đường đi ngắn nhất phụ thuộc thời gian.} Cooke và Halsey~\cite{cooke1966} là những người đầu tiên đặt bài toán
đường đi với thời gian di chuyển phụ thuộc thời điểm khởi hành, giải bằng quy hoạch động trên thời gian rời rạc.
Dreyfus~\cite{dreyfus1969} nhận xét rằng có thể mở rộng trực tiếp Dijkstra bằng cách gắn nhãn \emph{thời điểm đến} thay
cho khoảng cách. Kaufman và Smith~\cite{kaufman1993} chứng minh cách làm này cho kết quả đúng khi mạng thoả tính FIFO,
trong khi Orda và Rom~\cite{orda1990} phân tích trường hợp tổng quát, chỉ ra vai trò của việc cho phép hay cấm chờ tại
nút. Ichoua, Gendreau và Potvin~\cite{ichoua2003} đề xuất mô hình vận tốc thay đổi theo khoảng thời gian, trong đó thời
gian qua cạnh được tính bằng cách tích phân vận tốc, nhờ đó FIFO được bảo đảm theo cấu trúc. Delling và
Wagner~\cite{delling2009} tổng quan các kỹ thuật định tuyến phụ thuộc thời gian, bao gồm TD-A* và các biến thể có tiền xử lý.

\textbf{Tái lập kế hoạch trong môi trường thay đổi.} Trong robot học, D*~\cite{stentz1994} và D* Lite~\cite{koenig2002}
giải bài toán tìm đường khi chi phí cạnh chỉ được phát hiện dần trong quá trình di chuyển. Thay vì tìm lại từ đầu, D*
Lite duy trì lời giải tìm kiếm ngược từ đích và chỉ cập nhật phần bị ảnh hưởng khi một số chi phí cạnh thay đổi, nhờ đó
tiết kiệm đáng kể khi thay đổi mang tính cục bộ.

\textbf{Heuristic tham lam.} Một chiến lược đơn giản không cần mô hình toàn cục là luôn chọn cạnh ra hướng về đích
nhất (theo góc hoặc theo khoảng cách chim bay giảm nhiều nhất). Chiến lược này nhanh nhưng dễ đi vào ngõ cụt; chúng tôi
dùng nó làm mốc so sánh dưới cho các chính sách RL có cùng thông tin quan sát.

\subsection{Các hướng tiếp cận hiện đại}\label{sec:related:modern}

\textbf{Học tăng cường cho định tuyến.} Boyan và Littman~\cite{boyan1994} đề xuất Q-routing --- một trong những ứng dụng
sớm nhất của RL cho định tuyến gói tin trong mạng có tải thay đổi. Q-learning dạng bảng~\cite{watkins1992} học hàm giá
trị hành động bằng cập nhật thời gian--sai phân. Khi không gian trạng thái lớn, RL sâu dùng mạng nơ-ron xấp xỉ hàm giá
trị (DQN~\cite{mnih2015}) hoặc trực tiếp tối ưu chính sách (PPO~\cite{schulman2017}); các cài đặt chuẩn có trong thư viện
Stable-Baselines3~\cite{raffin2021}. Khi tập hành động hợp lệ thay đổi theo trạng thái (như số cạnh ra khác nhau ở mỗi
giao lộ), \emph{mặt nạ hành động} giúp gradient chính sách không lãng phí vào hành động không hợp lệ~\cite{huang2022}.
Định hình thưởng theo thế năng~\cite{ng1999} giúp giảm độ thưa của phần thưởng mà không làm thay đổi chính sách tối ưu;
học theo chương trình (curriculum learning)~\cite{bengio2009} bắt đầu từ nhiệm vụ dễ rồi tăng dần độ khó. Với quan sát
một phần, mạng hồi quy LSTM~\cite{hochreiter1997lstm} cho phép chính sách có bộ nhớ. Học bắt chước (behavior cloning) học
chính sách từ quỹ đạo của chuyên gia; DAgger~\cite{ross2011dagger} khắc phục hiện tượng lệch phân phối bằng cách thu
thập nhãn trên chính các trạng thái mà chính sách học được ghé thăm. Nền tảng chung của RL có trong~\cite{sutton2018}.

\textbf{Dự báo giao thông trên đồ thị.} Mạng tích chập đồ thị (GCN)~\cite{kipf2017gcn} tổng quát hoá phép tích chập lên
dữ liệu có cấu trúc đồ thị bằng cách trộn đặc trưng của mỗi nút với các nút kề. DCRNN~\cite{li2018dcrnn} mô hình hoá giao
thông như một quá trình khuếch tán trên đồ thị và kết hợp với mạng hồi quy; STGCN~\cite{yu2018stgcn} dùng các khối tích
chập không gian--thời gian. Các mô hình này dự báo vận tốc trên mạng cảm biến; trong khung của đề tài, một mô hình như
vậy có thể thay thế dự báo tĩnh $V_{\text{base}}$ làm đầu vào cho TD-A*.

\textbf{Định tuyến tính đến rủi ro và phân bổ giao thông.} Khi thời gian đi là ngẫu nhiên, người lái thường quan tâm đến
độ tin cậy (ví dụ thời gian đi ở phân vị 90\%) hơn là giá trị trung bình. Khi nhiều người cùng chọn tuyến, lý thuyết phân
bổ giao thông phân biệt cân bằng người dùng và tối ưu hệ thống theo nguyên lý Wardrop~\cite{wardrop1952}, với hàm trễ phổ
biến là hàm BPR~\cite{bpr1964}.

\textbf{Công cụ.} OSMnx~\cite{boeing2017} cho phép tải, làm sạch và phân tích mạng đường từ OpenStreetMap; Numba~\cite{lam2015numba}
biên dịch JIT mã Python số sang mã máy qua LLVM, giúp các vòng lặp tìm kiếm chạy nhanh gần bằng C.

\subsection{Phân tích điểm mạnh / hạn chế}\label{sec:related:compare}

Bảng~\ref{tab:compare} tóm tắt so sánh định tính giữa các họ phương pháp, dựa trên tài liệu và bản chất thuật toán.

\begin{table}[H]
\centering\small
\renewcommand{\arraystretch}{1.25}
\caption{So sánh định tính các họ phương pháp tìm đường.}
\label{tab:compare}
\begin{tabularx}{\linewidth}{>{\raggedright\arraybackslash}p{3.3cm}XX}
\toprule
\textbf{Họ phương pháp} & \textbf{Điểm mạnh} & \textbf{Hạn chế} \\
\midrule
Dijkstra / A* tĩnh & Đơn giản, tối ưu với trọng số cố định, rất nhanh & Bỏ qua biến thiên vận tốc theo giờ \\
TD-Dijkstra / TD-A* & Tối ưu khi mạng thoả FIFO, thời gian đa thức & Cần mô hình vận tốc tương lai; nhạy với dự báo sai \\
ALT, Contraction Hierarchies & Truy vấn rất nhanh trên mạng lớn & Cần tiền xử lý; bản phụ thuộc thời gian phức tạp \\
Replan, D* Lite & Thích nghi với quan sát mới trong khi đi & Chi phí tìm kiếm lặp; D* Lite cần chi phí tĩnh giữa các lần cập nhật \\
Dự báo GNN & Lan truyền thông tin từ cảm biến sang cạnh chưa quan sát & Cần dữ liệu huấn luyện; chỉ có ích khi nhiễu có tương quan \\
Định tuyến theo rủi ro & Giảm đuôi phân phối thời gian đi & Mục tiêu không cộng tính; tốn nhiều mẫu \\
RL chọn từng cạnh & Không cần mô hình chuyển trạng thái tường minh & Khám phá khó, khái quát kém, không bảo đảm tới đích \\
RL kết hợp tìm kiếm & Giữ bảo đảm tới đích và chất lượng của tìm kiếm & Không gian hành động hẹp nên lợi ích có thể nhỏ \\
\bottomrule
\end{tabularx}
\end{table}

\textbf{Khoảng trống nghiên cứu.} Nhiều công trình đánh giá RL cho định tuyến trên lưới tổng hợp, hoặc không so sánh với
TD-A* kết hợp tái lập kế hoạch trong cùng một môi trường; các kết quả vì vậy khó so sánh và dễ đánh giá quá cao vai trò
của RL. Đề tài hướng tới lấp khoảng trống đó ở quy mô một quận của Việt Nam: (i) dùng mạng đường thật từ OpenStreetMap;
(ii) so sánh mọi phương pháp trên \emph{cùng} tập kiểm tra, \emph{cùng} môi trường thực tái lập theo seed, \emph{cùng}
chỉ số và kiểm định thống kê ghép cặp; (iii) phân tích có hệ thống các cải tiến RL phổ biến và trình bày trung thực cả
những kết quả âm.
